\section{Frequentist and Bayesian Inference}
\label{sec:3.1}

Extracting physical information from noisy gamma-ray data requires a well-defined statistical language.
The analyses presented in this thesis draw on two complementary paradigms---frequentist and Bayesian inference---each providing distinct tools for parameter estimation, uncertainty quantification, and hypothesis testing.
This section introduces both frameworks from the shared foundation of the likelihood function, develops the key concepts used in the research papers, and discusses when each paradigm is most appropriate (according to personal bias).

%% --- 3.1.1 ---

\subsection{The Inference Problem}
\label{sec:3.1.1}

The starting point for any quantitative analysis in physics is the problem of inference: given a set of observed data $\mathbf{d}$, and a theoretical model $\mathcal{M}$ governed by parameters $\boldsymbol{\theta}$, what can we learn about $\boldsymbol{\theta}$ from $\mathbf{d}$?
In gamma-ray astrophysics, $\mathbf{d}$ might be a photon-count map recorded by the Fermi LAT, while $\boldsymbol{\theta}$ encodes the physical quantities of interest: source fluxes, spectral indices, or the dark matter annihilation cross-section $\langle \sigma v \rangle$.
The model $\mathcal{M}$ connects these parameters to the expected data through a forward process that may include astrophysical emission templates, instrumental response functions, and Poisson sampling of photon counts.

The connection between parameters and data is formalized by the likelihood function $\mathcal{L}(\boldsymbol{\theta}) = p(\mathbf{d} \,|\, \boldsymbol{\theta})$, which quantifies how well a given parameter choice explains the observed data \cite{Cranmer:2020}.
This function is the shared foundation of both major inference paradigms.
In the frequentist approach, the parameters $\boldsymbol{\theta}$ are treated as fixed but unknown quantities, and inference proceeds by identifying the parameter values that maximize $\mathcal{L}$.
In the Bayesian approach, $\boldsymbol{\theta}$ are themselves treated as random variables with a prior distribution $p(\boldsymbol{\theta})$, and the goal is to compute the posterior $p(\boldsymbol{\theta} \,|\, \mathbf{d})$ via Bayes' theorem \cite{bishop,Hastie:2009}.

These two philosophies are not in competition; they are complementary tools, each with distinct strengths suited to different aspects of the inference problem.
In the analyses presented in this thesis, both are used---sometimes on the same dataset---to extract physical information and quantify uncertainty.
The following two subsections develop each approach, starting from a common notation and building toward the specific formulations employed in the research papers.

%% --- 3.1.2 ---

\subsection{Frequentist Inference}
\label{sec:3.1.2}

In the frequentist framework, the model parameters $\boldsymbol{\theta}$ are fixed but unknown.
The data $\mathbf{d}$ are treated as one realization of a random process, and all statistical statements are made in terms of the long-run frequency properties of that process.
The central object is the likelihood function, which quantifies the probability of the observed data as a function of the model parameters \cite{Hastie:2009}.

\paragraph{Maximum Likelihood Estimation.}

The most common frequentist estimator is the maximum likelihood estimator (MLE), which identifies the parameter values that maximize the likelihood:
\begin{equation}
	\hat{\boldsymbol{\theta}}_{\mathrm{MLE}} = \arg\max_{\boldsymbol{\theta}} \, \mathcal{L}(\boldsymbol{\theta}).
	\label{eq:mle}
\end{equation}
As a concrete example, consider counting $N$ photons from a gamma-ray source.
If the counts follow a Poisson distribution with expected rate $\mu$, the likelihood is $\mathcal{L}(\mu) = e^{-\mu} \mu^N / N!
$, and the MLE is simply $\hat{\mu} = N$. While this example is trivial, the Poisson likelihood extends naturally to the pixel-by-pixel analysis of gamma-ray sky maps, where the joint likelihood is the product of Poisson terms over all pixels:
\begin{equation}
	\mathcal{L}(\boldsymbol{\theta}) = \prod_{a=1}^{N_{\mathrm{pix}}} \frac{\lambda_a(\boldsymbol{\theta})^{k_a} \, e^{-\lambda_a(\boldsymbol{\theta})}}{k_a!},
	\label{eq:poisson_likelihood}
\end{equation}
with $k_a$ the observed counts and $\lambda_a(\boldsymbol{\theta})$ the model prediction in pixel $a$.
This product-of-Poisson construction is the standard likelihood used in Fermi-LAT analyses \cite{Mattox:1996zz}.

\paragraph{Nuisance parameters and the profile likelihood.}

In realistic analyses, the parameter vector $\boldsymbol{\theta}$ contains both the parameters of interest $\boldsymbol{\psi}$ and a potentially large set of nuisance parameters $\boldsymbol{\eta}$ (background normalizations, instrumental systematics, source model parameters).
The profile likelihood isolates the dependence on $\boldsymbol{\psi}$ by maximizing over the nuisance parameters at each fixed value of $\boldsymbol{\psi}$:
\begin{equation}
	\mathcal{L}_{\mathrm{prof}}(\boldsymbol{\psi}) = \max_{\boldsymbol{\eta}} \, \mathcal{L}(\boldsymbol{\psi}, \boldsymbol{\eta}).
	\label{eq:profile_likelihood}
\end{equation}
This construction preserves the frequentist interpretation of the likelihood without requiring any prior distribution on $\boldsymbol{\eta}$.
An alternative approach is to integrate over the nuisance parameters within the likelihood framework itself, producing an integrated (or marginal) likelihood:
\begin{equation}
	\mathcal{L}_{\mathrm{int}}(\boldsymbol{\psi}) = \int \mathcal{L}(\boldsymbol{\psi}, \boldsymbol{\eta}) \, g(\boldsymbol{\eta}) \, d\boldsymbol{\eta},
	\label{eq:integrated_likelihood}
\end{equation}
where $g(\boldsymbol{\eta})$ is a weight function derived from independent auxiliary measurements or physical constraints on the nuisance parameters.
Although this resembles Bayesian marginalization in form, the integrated likelihood is interpreted differently: it is a tool for removing $\boldsymbol{\eta}$ from the likelihood rather than a posterior probability statement about $\boldsymbol{\psi}$.

\paragraph{Hypothesis testing and confidence intervals.}

To compare competing hypotheses---for instance, whether a gamma-ray source exists at a given position---we construct a likelihood-ratio test statistic (TS):
\begin{equation}
	\mathrm{TS} = -2 \ln \frac{\mathcal{L}_0}{\mathcal{L}_1},
	\label{eq:ts}
\end{equation}
where $\mathcal{L}_0$ is the likelihood under the null hypothesis and $\mathcal{L}_1$ the likelihood under the alternative (or, more commonly, the globally maximized likelihood) \cite{Mattox:1996zz}.
By Wilks' theorem, under mild regularity conditions the TS is asymptotically distributed as a $\chi^2$ random variable with degrees of freedom equal to the difference in free parameters between the two hypotheses \cite{Wilks:1938}, which lets us convert a TS value into a statistical significance without resorting to explicit Monte Carlo simulations.
The same asymptotic result underpins the construction of confidence intervals, understood as the inverted test: rather than fixing a single threshold and asking whether the null is rejected, we scan the parameter of interest and retain the range over which the profile-likelihood TS stays below the appropriate $\chi^2$ quantile.
For a 95\% confidence-level upper bound on a single parameter, the conventional choice is to identify the value at which $\mathrm{TS} = 3.84$---the 95\% quantile of $\chi^2_1$, which is conservative for a one-sided bound.
Hypothesis testing and interval estimation are thus two uses of the same construction: the same test statistic and the same quantile quantify a detection significance and delimit a confidence interval.

\paragraph{Bootstrap confidence intervals.}

The $\chi^2$ calibration of the TS rests on the regularity conditions of Wilks' theorem, and these can fail in situations that recur in gamma-ray analyses: when a parameter lies on a physical boundary---a flux or a cross-section that cannot be negative---or when the sample is small enough that the asymptotic regime has not yet set in.
In such cases the bootstrap provides a general-purpose remedy \cite{Efron:1979}: we repeatedly refit the estimator, either on datasets resampled from the observed data (the \emph{nonparametric} bootstrap) or on synthetic datasets drawn from the fitted model (the \emph{parametric} bootstrap), and read the confidence interval directly from the spread of the re-fitted estimates---the containment interval of their empirical distribution.
In Chapter~\ref{ch:5}, the uncertainty band on the dark matter annihilation cross-section bounds is obtained in exactly this way: for each dark matter mass we sample 100 synthetic datasets from the best-fit mixture model, re-derive the bound on $\langle \sigma v \rangle$ on each of them, and take the 95\% containment interval of the resulting distribution of bounds as the band, with its upper edge further adjusted to conservatively account for systematics associated with the presence of 4FGL sources.

\paragraph{Goodness of fit.}

A different question arises when we wish to assess the absolute compatibility of the data with a single model, without testing it against any specific alternative.
The Kolmogorov--Smirnov test \cite{KS_test} addresses this goodness-of-fit question by taking as its test statistic the supremum of the absolute difference between the empirical cumulative distribution function of the data and the cumulative distribution predicted by the model.
When the null model is fully specified, the distribution of this statistic does not depend on the underlying distribution, which makes it a convenient distribution-free diagnostic whenever no particular alternative hypothesis is in play.

\paragraph{Non-nested model comparison.}

The likelihood-ratio test presumes that the two hypotheses are nested, with one model recovered as a special case of the other; when this is not so, we need a criterion able to rank models of different structure.
Information criteria fill this gap.
The Akaike Information Criterion \cite{Akaike:1974} is defined as
\begin{equation}
	\mathrm{AIC} = -2 \ln \mathcal{L}_{\max} + 2k,
	\label{eq:aic}
\end{equation}
where $\mathcal{L}_{\max}$ is the maximum value of the likelihood and $k$ the number of free parameters.
The AIC provides an estimator of the relative Kullback--Leibler divergence between each candidate model and the data-generating process, with the $2k$ term penalising additional parameters and thus discouraging overfitting.
Model selection then proceeds by minimising the AIC across the candidates; \blue{differences in AIC will then determine a hierarchy of relative preference.}

\paragraph{Application preview.}

The frequentist framework developed above is used extensively in this thesis:
\begin{itemize}
	\item In Chapter~\ref{ch:4}, the MSP luminosity function is constrained via MLE applied to an integrated likelihood, and model comparison is performed through $-2\Delta\ln\mathcal{L}$ values.
	\item In Chapter~\ref{ch:5}, a profile likelihood test statistic is used with Wilks' theorem to set 95\% CL upper bounds on the dark matter annihilation cross-section, with the uncertainty band on those bounds derived from a parametric bootstrap; the AIC is additionally used to compare competing models for the unassociated source population.
	\item In Chapter~\ref{ch:7}, a two-sample KS test compares the cumulative pixel-wise TS distributions of the real \Fermi-LAT sky against synthetic realizations drawn from the source-count distribution, providing the statistical criterion that defines sub-threshold source candidates.
	\item In Chapter~\ref{ch:8}, a $\Delta\chi^2$ test statistic compares astrophysical-only and astrophysical-plus-DM hypotheses for the cross-correlation angular power spectrum.
\end{itemize}

%% --- 3.1.3 ---

\subsection{Bayesian Inference}
\label{sec:3.1.3}
The Bayesian framework treats the model parameters $\boldsymbol{\theta}$ not as fixed unknowns but as random variables associated with a probability distribution.
Before observing any data, our state of knowledge about $\boldsymbol{\theta}$ is encoded in a prior distribution $p(\boldsymbol{\theta})$.
After observing data $\mathbf{d}$, this knowledge is updated through Bayes' theorem:
\begin{equation}
	p(\boldsymbol{\theta} \,|\, \mathbf{d}) = \frac{p(\mathbf{d} \,|\, \boldsymbol{\theta}) \, p(\boldsymbol{\theta})}{p(\mathbf{d})},
	\label{eq:bayes}
\end{equation}
where $p(\mathbf{d} \,|\, \boldsymbol{\theta}) = \mathcal{L}(\boldsymbol{\theta})$ is the same likelihood function used in the frequentist approach, and $p(\mathbf{d}) = \int p(\mathbf{d} \,|\, \boldsymbol{\theta}') \, p(\boldsymbol{\theta}') \, d\boldsymbol{\theta}'$ is the evidence, a normalization constant that integrates the likelihood weighted by the prior over the full parameter space \cite{bishop}.
Within the assumed model, the posterior $p(\boldsymbol{\theta} \,|\, \mathbf{d})$ is the complete answer to the inference problem: it encodes not only the best-fit values but also the full uncertainty structure, including correlations between parameters.

\paragraph{Prior, posterior, and evidence.}

The choice of prior $p(\boldsymbol{\theta})$ is one of the defining features of the Bayesian approach.
Informative priors can incorporate physical constraints---for example, requiring fluxes to be non-negative or restricting spectral indices to physically motivated ranges---and can improve inference when data are sparse.
Uninformative or weakly informative priors are used when prior knowledge is limited, though the choice is never fully objective \cite{Hastie:2009}.
For complex models, the evidence $p(\mathbf{d})$ is typically intractable to compute, as it requires integration over the full parameter space, but it plays an important role in Bayesian model comparison, where the ratio of evidences between competing models (the Bayes factor) provides a natural criterion for model selection \cite{Kass:1995}.

\paragraph{Credible intervals versus confidence intervals.}

In the Bayesian framework, uncertainty is quantified through credible intervals: a 95\% credible interval is a region of parameter space containing 95\% of the posterior probability.
This contrasts with the frequentist confidence interval, which guarantees that 95\% of intervals constructed from repeated experiments will contain the true parameter value.
The distinction is subtle in practice---for well-constrained problems, the two often agree numerically---but the conceptual difference matters: a credible interval makes a direct probability statement about the parameter given the data, while a confidence interval makes a statement about the procedure.

\paragraph{Marginalization over nuisance parameters.}

As in the frequentist case, realistic Bayesian analyses involve nuisance parameters.
The Bayesian approach handles them by marginalizing the joint posterior:
\begin{equation}
	p(\boldsymbol{\psi} \,|\, \mathbf{d}) = \int p(\boldsymbol{\psi}, \boldsymbol{\eta} \,|\, \mathbf{d}) \, d\boldsymbol{\eta},
	\label{eq:bayesian_marginalization}
\end{equation}
which integrates out the nuisance parameters $\boldsymbol{\eta}$ to yield the marginal posterior for the parameters of interest $\boldsymbol{\psi}$.
This is conceptually analogous to the frequentist integrated likelihood (see Section~\ref{sec:3.1.2}), but the probabilistic interpretation differs: the Bayesian marginal posterior is a fully coherent probability distribution over $\boldsymbol{\psi}$, conditioned on both the data and the prior.

When the likelihood is intractable---as is often the case for complex forward models involving high-dimensional simulations---neural networks offer a route to Bayesian inference through implicit likelihood methods, in which a density estimator learns the posterior directly from simulated data--parameter pairs.
These approaches are developed in Section~\ref{sec:3.2}.

\paragraph{Application preview.}

In Chapter~\ref{ch:6}, the Bayesian framework is implemented via simulation-based inference to reconstruct the source-count distribution $dN/dS$ from gamma-ray sky maps, with posterior uncertainty estimates validated against an independent frequentist cross-check (Section~\ref{sec:3.1.4}).

%% --- 3.1.4 ---

\subsection{When to Use Which}
\label{sec:3.1.4}

For most inference tasks encountered in practice---parameter estimation, upper-limit setting, model comparison---both the frequentist and Bayesian frameworks can reach the same scientific conclusion.
The choice between them is usually pragmatic rather than fundamental, driven by the scientific question at hand, computational considerations, and the conventions of the relevant research community.
Community practice already carries considerable weight: high-energy particle physics and the Fermi-LAT collaboration have traditionally favored frequentist methods---source detection via the TS, upper limits via the profile likelihood and Wilks' theorem, catalog pipelines built on maximum-likelihood fitting \cite{Mattox:1996zz}---while cosmology and large-scale-structure analyses have largely adopted Bayesian Markov Chain Monte Carlo (MCMC) methods \cite{Hastie:2009}.
Neither convention reflects a judgment on the foundations of probability; both reflect what works well for the typical dimensionality and structure of the problems each field encounters most often.

Beyond community defaults, the central question is whether the analysis requires a full characterization of uncertainty or whether a point estimate with a simple error bar suffices.
When the posterior structure matters---correlations between parameters, degeneracies, multi-modality, or the marginalization of many nuisance parameters---the Bayesian framework provides this information as its default output (Section~\ref{sec:3.1.3}).
In high-dimensional parameter spaces, however, MCMC itself becomes expensive.
Frequentist point estimation can remain tractable in the same regime through gradient-based optimization or iterative algorithms such as expectation--maximization (EM) \blue{(Chapter~\ref{ch:5})}, which converge efficiently without exploring the full parameter space.
What is given up is the posterior structure itself, not uncertainty quantification: confidence intervals remain available through the bootstrap (Section~\ref{sec:3.1.2}), at the cost of repeated refitting.
When a one-dimensional confidence interval or upper limit is the target, as in standard source detection or cross-section bounds, the frequentist route is typically more direct.

A second discriminating factor is the role of external information.
Bayesian inference provides a coherent mechanism for incorporating physical constraints---positivity of fluxes, calibration ranges from auxiliary measurements, results from complementary experiments---through the prior.
Frequentist methods are preferred when the goal is to present a result with minimal sensitivity to external assumptions, particularly in the data-sparse regime where the prior can considerably shape the posterior.
This prior dependence extends to model comparison: the Bayes factor offers a principled comparison between arbitrary model pairs but is sensitive to the prior volume, especially when priors are vague or the models differ substantially in complexity \cite{Kass:1995}; the frequentist likelihood-ratio test (Section~\ref{sec:3.1.2}) is prior-free but restricted to nested hypotheses; information criteria such as the AIC apply more broadly but produce only relative rankings without significance levels.

In the limit of abundant, constraining data, these distinctions fade: the likelihood dominates any reasonable prior, and Bayesian credible intervals coincide numerically with frequentist confidence intervals.
The choice becomes consequential precisely when data are sparse, priors are informative, or the parameter space is high-dimensional.
In the analyses presented in this thesis, both paradigms are used, and in some cases they are applied to the same problem as independent cross-checks.
In Chapter~\ref{ch:6}, a simulation-based inference approach provides Bayesian posterior estimates on the reconstructed $dN/dS$, which are validated against a frequentist analysis of the same data where confidence intervals are extracted from the empirical distribution of residuals across the validation set.
The frequentist cross-check also provides a formal estimate of the model's bias, which is found to be negligible across the flux range of interest.
The compatibility of the two independent error estimates provides confidence that the uncertainty quantification is robust.
